% =====================================================================
%  CAPÍTULO 1 -- INTRODUÇÃO
%  Formato: PROJETO APLICADO (relatório técnico + protótipo)
%
%  O Manual recomenda texto corrido, sem dividir a introdução em seções.
% =====================================================================

\chapter{Introdução e objetivos}
\label{cap:introducao}

% (a) Contexto e problema
A prática regular de lutas e artes marciais está associada a benefícios
que vão além do condicionamento físico. Em revisão sistemática com
meta-análise, \citeonline{moore2020} encontraram efeitos positivos do
treinamento em artes marciais sobre o bem-estar e a saúde mental dos
praticantes, incluindo a redução de sintomas de ansiedade e depressão.
Esses benefícios, porém, só aparecem com a continuidade do treino. Na
vivência do autor como praticante, é comum que alunos abandonem a
atividade antes de consolidar o hábito, seja pela percepção lenta da
própria evolução, pela falta de metas claras, pelo distanciamento do
professor ou por mudanças na rotina; o levantamento com praticantes e
ex-praticantes descrito no \autoref{cap:testes} procurou, justamente,
confirmar ou refutar essas impressões. Para o aluno, o abandono interrompe
o ciclo em que os ganhos se tornariam perceptíveis; para a academia e o
professor, representa perda de receita e de vínculo com a turma.

% (b) Solução proposta
Para enfrentar esse problema, este projeto propõe o \textbf{RoundUp}, um
aplicativo móvel que aplica \textbf{gamificação} -- o uso de elementos de
design de jogos em contextos que não são jogos \cite{deterding2011} -- ao
acompanhamento dos treinos. Para o aluno, o aplicativo oferece mecânicas
como pontos de experiência, sequências de treino, medalhas e uma barra de
progresso rumo à próxima graduação. Para o professor, oferece apoio à
gestão: registro de presença, criação de missões para a turma e alertas
sobre alunos com queda de frequência. O sistema distingue três perfis de
usuário: o \textbf{aluno}, o \textbf{professor}, que pode atuar de forma
autônoma, e a \textbf{academia}, que reúne professores, turmas e alunos.
O protótipo resultante é de natureza \textbf{digital} (Regulamento,
Art.~10): um aplicativo funcional em React Native com \textit{back-end}
no Firebase.

% Natureza de ideação do projeto
Por se tratar de um domínio em que o valor de cada mecânica de jogo para
o público ainda é pouco conhecido, o projeto foi conduzido como um
\textbf{projeto de ideação} do autor. A partir da própria vivência, da
literatura e de aplicativos de referência, o autor elaborou ideias de
funcionalidades e as converteu em \textbf{requisitos iniciais},
consolidados em setembro de 2026 e usados desde então para construir o
aplicativo. Essas ideias foram tratadas como \textbf{hipóteses}: em vez
de serem tomadas como certas, foram submetidas a um questionário
(\textit{survey}) que serviu para \textbf{confirmar ou refutar} cada
uma e para \textbf{revelar ideias novas}, não pensadas no início. O
questionário foi respondido por praticantes e ex-praticantes reais e,
para cobrir os perfis ausentes dessa coleta, como o de professor, também
por \textbf{personas sintéticas}: 12 personas, elaboradas com apoio de
um modelo de linguagem de grande porte (LLM), orientaram um
\textit{script} que gerou 65 \textbf{respostas sintéticas}. Essas
respostas trouxeram ideias novas e um sinal complementar sobre cada
requisito, mas a decisão sobre cada um se apoiou nos dados reais; elas
serviram ainda para testar o instrumento e o procedimento de análise. A
\autoref{fig:ciclo-ideacao}, no \autoref{cap:requisitos}, ilustra esse
ciclo.

A construção do aplicativo também usou inteligência artificial, desta vez
como \textbf{ferramenta de desenvolvimento}: um agente de codificação
executou boa parte da escrita do código a partir de instruções,
requisitos e planos definidos pelo autor, que concentrou o seu trabalho no
planejamento, no projeto da solução e no controle de qualidade, por meio
de testes automatizados, de regras de segurança verificadas no servidor e
de execução em aparelho (\autoref{sec:dev-ia}).

% (c) Objetivos
O objetivo geral deste projeto é \textbf{idear, desenvolver e verificar
o RoundUp, um aplicativo gamificado para academias de lutas e artes
marciais, voltado a aumentar o engajamento e a retenção dos alunos e a
apoiar o professor na gestão da turma, avaliando com potenciais usuários
as ideias que orientam os seus requisitos}. Os objetivos específicos são:

\begin{enumerate}
  \item identificar fatores de motivação e de evasão de praticantes de
        lutas e artes marciais;
  \item elaborar ideias de mecânicas de gamificação e de apoio ao
        professor e convertê-las em requisitos iniciais rastreáveis;
  \item confirmar ou refutar essas ideias com usuários reais, por meio
        de um questionário, e identificar ideias novas a partir das
        respostas;
  \item construir personas sintéticas e usá-las para gerar ideias não
        previstas, como sinal complementar sobre os requisitos e como
        dados de teste do questionário e do procedimento de análise, sem
        misturá-las aos dados reais;
  \item modelar e implementar um protótipo do aplicativo com os
        requisitos priorizados, empregando agentes de IA como ferramenta
        de desenvolvimento sob o planejamento, o projeto e o controle de
        qualidade do autor;
  \item considerar, desde a concepção, a conformidade com a Lei Geral de
        Proteção de Dados Pessoais (LGPD), inclusive no tratamento de dados
        de menores de idade \cite{lgpd2018}.
\end{enumerate}

% (d) Justificativa
O projeto tem relevância social, por contribuir para a permanência de
alunos em uma atividade associada à saúde física e mental; relevância
técnica, por articular engenharia de requisitos, gamificação,
desenvolvimento móvel e proteção de dados pessoais em uma solução
prática; e relevância metodológica em duas frentes. A primeira é
documentar o uso de personas e respostas sintéticas como fonte de ideias,
sinal complementar e dados de teste de um instrumento de pesquisa,
deixando explícito o que elas podem e o que não podem mostrar. A segunda
é registrar, em um projeto real, como o desenvolvimento com agentes de IA
desloca o papel do desenvolvedor. O uso de IA já é a regra na profissão:
cerca de 90\% dos profissionais de desenvolvimento entrevistados em 2025
relatam usá-la no trabalho \cite{dora2025}, e 84\% dos respondentes de
outra pesquisa internacional usam ou pretendem usar essas ferramentas
\cite{stackoverflow2025}. Os ganhos, porém, não são automáticos:
estudos controlados encontraram desde tarefas concluídas 55,8\% mais
rápido \cite{peng2023} até desenvolvedores experientes 19\% mais
lentos com IA, embora acreditassem estar mais rápidos \cite{becker2025}; e
a confiança no código gerado segue baixa \cite{dora2025,stackoverflow2025}.
Nesse cenário, o que diferencia o resultado é a capacidade do
desenvolvedor de planejar, especificar, projetar e verificar o que o
agente produz, e é essa prática que este relatório documenta. Soma-se a isso a vivência do autor
com as artes marciais, que facilita o acesso a praticantes e professores.

% (e) Escopo
O escopo do projeto concentra-se em academias de lutas e artes marciais
no Brasil e em duas perspectivas de uso: a do aluno (engajamento) e a do
professor (gestão). Ficam fora do escopo: outras modalidades esportivas,
prescrição de treinos ou orientação nutricional, integração com
dispositivos vestíveis, análise de movimentos por visão computacional e
aspectos comerciais, como monetização e escalabilidade em produção. A
avaliação do impacto de longo prazo sobre a evasão também ultrapassa o
prazo do TCC.

O recorte abrange lutas e artes marciais em geral: o aplicativo não depende de
uma modalidade específica, e o usuário escolhe a sua no cadastro.

% (f) Estrutura do relatório
Este relatório está organizado da seguinte forma: o
\autoref{cap:requisitos} apresenta os fundamentos que orientaram a
ideação, o processo de ideação e os requisitos do projeto; o
\autoref{cap:modelagem}, as modelagens funcional, estática e dinâmica; o
\autoref{cap:implementacao}, as estratégias de implementação do
protótipo, incluindo o desenvolvimento assistido por agentes de IA; o \autoref{cap:testes}, a avaliação das
ideias com usuários reais, o uso das respostas sintéticas na ideação e no
teste do questionário, e os testes do protótipo; e o \autoref{cap:resultados}, os resultados e as
considerações finais.